  % =====================================================================
  %  IEEE-STYLE CONFERENCE REPORT
  %  MLE-AI-201 Machine Learning -- Final Project
  % =====================================================================
  \documentclass[10pt,twocolumn,a4paper]{article}

  % ---- Geometry: IEEE-like margins / two narrow columns ----------------
  \usepackage[a4paper,top=0.75in,bottom=1.0in,left=0.68in,right=0.68in,columnsep=0.25in]{geometry}

  % ---- Times-like font (standard, ships with every TeX distro) ---------
  \usepackage{mathptmx}          % Times text + math
  \usepackage[T1]{fontenc}
  \usepackage[utf8]{inputenc}    % allows accented / Azerbaijani chars
  \DeclareUnicodeCharacter{0259}{e}
  \DeclareUnicodeCharacter{018F}{E}

  % ---- Common, dependency-free packages --------------------------------
  \usepackage{graphicx}
  \usepackage{amsmath,amssymb}
  \usepackage{booktabs}
  \usepackage{array}
  \usepackage{titlesec}
  \usepackage{caption}
  \usepackage{cite}
  \usepackage[hidelinks]{hyperref}
  \usepackage{url}

  % ---- IEEE-style section headings -------------------------------------
  \renewcommand{\thesection}{\Roman{section}}
  \renewcommand{\thesubsection}{\Alph{subsection}}
  \titleformat{\section}[block]{\centering\normalsize\scshape}{\thesection.}{0.5em}{}
  \titleformat{\subsection}[block]{\normalsize\itshape}{\thesubsection.}{0.5em}{}
  \titlespacing*{\section}{0pt}{1.4ex plus 1ex minus .2ex}{1ex}
  \titlespacing*{\subsection}{0pt}{1.0ex plus 1ex minus .2ex}{0.6ex}

  % ---- Abstract / Index-Terms in the IEEE boldface-dash style ----------
  \newcommand{\abstractkw}[1]{\noindent\textbf{\textit{#1}}}

  \captionsetup{font=small,labelsep=period}
  \setlength{\parindent}{1em}

  % =====================================================================
  \begin{document}

  % ------------------------- TITLE BLOCK --------------------------------
  \twocolumn[
    \begin{@twocolumnfalse}
    \begin{center}
      {\LARGE\bfseries Predicting Real-Estate Prices and Price Tiers from Scratch:\\
      \Large Decision Tree and Pegasos SVM Tested on bina.az\par}
      \vspace{1.2em}
      % Authors: Surname, Given name --- alphabetical by surname.
      {\large
        Alasgarov Teymur\quad Heydarov Tofiq\quad Salahov Parvin\par}
      \vspace{0.5em}
      {\itshape MLE-AI-201 Machine Learning, AI Academy --- Cohort II\par}
      \vspace{1.4em}
    \end{center}
    \begin{quote}
      \small
      \abstractkw{Abstract---}%
      Estimating property prices is hard in a market like Azerbaijan's. Prices are very skewed, location matters a lot, and listings scraped from websites are messy. In this project we built two models completely from scratch using only NumPy. The first is a decision tree (it uses MSE for regression, and Gini impurity or entropy for classification). The second is a Support Vector Machine trained with the Pegasos algorithm, with Random Fourier Features (RFF) added so it can learn non-linear patterns.

      We tested both models on 93,150 cleaned listings from \texttt{bina.az}, with two tasks: predicting the log-price (Task A), and classifying whether a property is above or below the median price (Task B).

      On Task A, our tree got an RMSE of 0.5036 ($R^2 = 0.5992$). scikit-learn got 0.5026 ($R^2 = 0.6007$), so the difference is only 0.0010. On Task B, our tree got $90.79\%$ test accuracy ($F_1 = 0.9131$, ROC-AUC 0.9675), which matches scikit-learn ($F_1 = 0.9131$, ROC-AUC 0.9674). Our Pegasos SVM with RFF reached $84.39\%$ accuracy ($F_1 = 0.8512$, ROC-AUC 0.9271), and its ROC-AUC was higher than the standard linear SGD baseline (0.9198). 5-fold cross-validation shows that the results are stable ($F_1 = 0.9123 \pm 0.0016$).

      Our code, tuning experiments, and scripts are available at \href{https://github.com/ParvinSalahov/ML_Final_Project}{\texttt{github.com/ParvinSalahov/ML\_Final\_Project}} (tag \texttt{v1.0-final}).
      
      \vspace{0.6em}
      \abstractkw{Index Terms---}Decision tree, Support Vector Machine, Pegasos, Random Fourier Features, real-estate valuation, reproducible benchmarking.
    \end{quote}
    \vspace{1.0em}
    \end{@twocolumnfalse}
  ]

  % =====================================================================
  \section{Introduction}
  Automated valuation models (AVMs) are used for mortgage decisions, market transparency, and risk management. But raw listings from \texttt{bina.az} are tricky to work with. Prices go from 11 AZN up to 627,130 AZN, the descriptions are written in free-form Azerbaijani with inconsistent wording, and many categories are very imbalanced.

  In this project we wanted to show that models we code ourselves from the basic math can give results close to professional libraries. Our main contributions are:
  \begin{itemize}
    \item \textbf{From-scratch algorithms}: A recursive binary decision tree (MSE, Gini, and entropy) and a hinge-loss SVM trained with Pegasos, using Random Fourier Features to get non-linear decision boundaries. Both are vectorised with NumPy.
    \item \textbf{Leakage-free features}: Careful data cleaning, extracting information from the text, frequency encoding, and computing every statistic only on the training data.
    \item \textbf{Benchmarking}: We compare against scikit-learn on 93,150 samples, with 5-fold cross-validation, residual analysis, and hyperparameter sweeps.
    \item \textbf{Reproducibility}: One command (\texttt{python -m src.run\_all}) recreates all 19 figures and tables.
  \end{itemize}

  % =====================================================================
  \section{Related Work / Background}
  Decision trees split the feature space with simple yes/no rules, choosing the split that makes the groups as pure as possible~\cite{breiman1984}. They are greedy, but they handle mixed numeric and categorical features, non-linear patterns, and different feature scales without needing standardization.

  A soft-margin SVM looks for the hyperplane with the widest margin while allowing a few violations (the hinge loss). Instead of solving the usual dual problem with SMO, Shalev-Shwartz et al.~\cite{pegasos} introduced Pegasos, which works directly on the primal objective with stochastic sub-gradient steps. It reaches an $\epsilon$-accurate solution in $\mathcal{O}(1/(\lambda\epsilon))$ iterations, no matter how large the dataset is. To capture non-linear patterns in the market, we map features with Random Fourier Features~\cite{rahimi2007}, which approximate the Gaussian RBF kernel $\mathcal{K}(\mathbf{x}, \mathbf{z}) = \exp(-\gamma \|\mathbf{x} - \mathbf{z}\|^2)$.

  % =====================================================================
  \section{Data and Preprocessing}
  \subsection{Dataset Cleaning}
  The raw dataset has 100,775 listings scraped from \texttt{bina.az}. We cleaned it in these steps:
  \begin{enumerate}
    \item \textbf{Removing duplicates}: Removing exact duplicate listings left 100,773 rows.
    \item \textbf{Extracting information}: Regular expressions pull the number of rooms and the renovation status (\textit{təmir}) out of the free-text description and attributes fields.
    \item \textbf{Filtering extreme rows}: We removed entries that make no physical sense (for example zero area, non-residential anomalies, and corrupted outliers). This left 93,150 valid records with 28 features.
    \item \textbf{Transforming the target}: Raw prices are strongly skewed (skewness $0.877$). We use $y = \log(1 + \text{price})$, which shrinks the huge price range and stabilizes the variance (skewness becomes $-1.340$).
    \item \textbf{Preventing leakage}: We dropped administrative IDs, duplicate URLs, and price-per-unit ratios, because they leak the target. Imputation values, frequency encodings, and scaling statistics were all computed on the training split only.
  \end{enumerate}

  \begin{figure}[ht]
    \centering
    \includegraphics[width=\linewidth]{figures/explore_price_distribution.png}
    \caption{Distribution of raw prices (left, strongly skewed) compared with log-transformed prices $\log(1 + \text{price})$ (right), on the training data.}
    \label{fig:price_dist}
  \end{figure}

  \begin{table}[ht]
    \centering
    \caption{Dataset summary and split statistics.}
    \label{tab:data}
    \resizebox{\columnwidth}{!}{%
    \begin{tabular}{@{}lrr@{}}
      \toprule
      Property & Absolute Count & Share \\
      \midrule
      Raw input rows              & 100,775 & 100.0\% \\
      After deduplication         & 100,773 & 100.0\% \\
      Total clean analytical rows & 93,150  & 100.0\% \\
      Train split                 & 62,922  & 67.5\% \\
      Validation split            & 15,114  & 16.2\% \\
      Test split                  & 15,114  & 16.2\% \\
      Engineered input features   & 28      & --- \\
      Task B train median price   & 200,000 AZN & $\log p = 12.21$ \\
      Task B balance (Tr/Val/Te)  & 49.4 / 54.1 / 54.2 & (\% Premium) \\
      \bottomrule
    \end{tabular}%
    }
  \end{table}

  \subsection{The Two Tasks}
  We formulate two complementary learning challenges:
  \begin{itemize}
    \item \textbf{Task A (regression)}: Predict the log-price $y = \log(1 + \text{price})$.
    \item \textbf{Task B (classification)}: Decide whether a property is premium (1) or standard (0), using the training median as the cutoff ($\tau = 12.2061$, about 200,000 AZN). The classes are fairly balanced; the test set has 54.2\% premium and 45.8\% standard.
  \end{itemize}

  % =====================================================================
  \section{Methodology}

  \subsection{Decision Tree (From Scratch)}
  The tree builds its splits greedily and recursively. For every feature $j$ and every candidate threshold $\theta$, we measure how much the split reduces impurity:
  \begin{equation}
    \Delta I(S, j, \theta) = I(S) - \frac{|S_L|}{|S|} I(S_L) - \frac{|S_R|}{|S|} I(S_R).
  \end{equation}
  Here $S_L = \{(\mathbf{x},y) \in S \mid x_j \le \theta\}$ is the group of samples with $x_j \le \theta$, and $S_R = S \setminus S_L$ is everything else. In other words, we compare the impurity before the split with the weighted impurity of the two groups after it.

  For regression (Task A), the impurity is the mean squared error around the node's mean:
  \begin{equation}
    I_{\text{MSE}}(S) = \frac{1}{|S|}\sum_{i \in S} (y_i - \bar{y}_S)^2, \quad \bar{y}_S = \frac{1}{|S|}\sum_{i \in S} y_i.
  \end{equation}
  For classification (Task B), we implemented two options, where $p_k = \frac{1}{|S|}\sum_{i \in S} \mathbb{I}(y_i = k)$ is the share of samples of class $k$ in the node:
  \begin{align}
    I_{\text{Gini}}(S) &= 1 - \sum_{k=0}^{1} p_k^2, \\
    I_{\text{Entropy}}(S) &= - \sum_{k=0}^{1} p_k \log_2(p_k + \epsilon).
  \end{align}
  A node stops splitting when any of these happens: the depth reaches $d_{\max}$, the node has fewer than $2 \cdot n_{\text{leaf}}$ samples, or a split would leave fewer than $n_{\text{leaf}}$ samples on either side. A leaf predicts the mean $\bar{y}$ for regression, and the class probabilities $[p_0, p_1]$ for classification.

  \subsection{Support Vector Machine with Pegasos (From Scratch)}
  For binary classification with labels $y_i \in \{-1, +1\}$, we minimize the regularized hinge-loss objective:
  \begin{equation}
    J(\mathbf{w}, b) = \frac{\lambda}{2} \|\mathbf{w}\|_2^2 + \frac{1}{n}\sum_{i=1}^{n} \max\big(0, 1 - y_i(\mathbf{w}^\top \mathbf{x}_i + b)\big).
  \end{equation}
  The value $\lambda > 0$ controls the regularization ($C \approx 1/(n\lambda)$). At step $t$ the learning rate is $\eta_t = \frac{1}{\lambda t}$, and we pick one random sample $(\mathbf{x}_t, y_t)$. If it violates the margin, meaning $y_t(\mathbf{w}_t^\top \mathbf{x}_t + b_t) < 1$, we update:
  \begin{align}
    \mathbf{w}_{t+1} &= (1 - \eta_t \lambda)\mathbf{w}_t + \eta_t y_t \mathbf{x}_t, \\
    b_{t+1} &= b_t + \eta_t y_t.
  \end{align}
  Otherwise we only shrink the weights: $\mathbf{w}_{t+1} = (1 - \eta_t \lambda)\mathbf{w}_t$.

  To get non-linear boundaries cheaply, we map the standardized features $\mathbf{x} \in \mathbb{R}^d$ into a $D$-dimensional random Fourier space $\mathbf{z}(\mathbf{x}) \in \mathbb{R}^D$:
  \begin{equation}
    z_j(\mathbf{x}) = \sqrt{\frac{2}{D}} \cos(\mathbf{\omega}_j^\top \mathbf{x} + \beta_j).
  \end{equation}
  The frequencies $\mathbf{\omega}_j$ are drawn from a normal distribution $\mathcal{N}(\mathbf{0}, 2\gamma \mathbf{I})$, and the phase shifts $\beta_j$ are drawn uniformly from $[0, 2\pi]$. This way we avoid building the huge $n \times n$ kernel matrix of the dual SVM.

  % =====================================================================
  \section{Experimental Setup}
  \subsection{Evaluation Metrics}
  All final scores come from the test split that the models never saw ($n = 15,114$):
  \begin{itemize}
    \item \textbf{Task A}: RMSE, MAE, and $R^2$.
    \item \textbf{Task B}: Accuracy, Precision, Recall, Macro $F_1$, ROC-AUC, and PR-AUC.
  \end{itemize}

  \subsection{Hyperparameter Tuning}
  We tuned hyperparameters only on the validation set ($n = 15,114$):
  \begin{itemize}
    \item \textbf{Decision tree}: We tried max depth from 1 to 18, both criteria (Gini and entropy), and leaf size from 1 to 100. The best setting was $\text{criterion} = \text{entropy}$, $d = 12$, $n_{\text{leaf}} = 10$.
    \item \textbf{Pegasos SVM}: We tried $\lambda$ from $10^{-5}$ to $1.0$ and the RFF scale $\gamma$ from $0.1$ to $10.0$, with 256 components and 30 epochs. The best setting was $\lambda = 10^{-4}$ ($C \approx 0.1589$) and $\gamma = 0.1$.
  \end{itemize}

  % =====================================================================
  \section{Results and Benchmarking}

  \begin{table*}[t]
    \centering
    \caption{Consolidated Headline Results on Unseen Test Split ($n=15,114$).}
    \label{tab:results}
    \resizebox{\textwidth}{!}{%
    \begin{tabular}{@{}lcccccccc@{}}
      \toprule
      & \multicolumn{3}{c}{\textbf{Task A: Log-Price Regression}} & \multicolumn{4}{c}{\textbf{Task B: Price-Tier Classification}} & \textbf{Compute} \\
      \cmidrule(lr){2-4} \cmidrule(lr){5-8} \cmidrule(lr){9-9}
      Model & RMSE $\downarrow$ & MAE $\downarrow$ & $R^2 \uparrow$ & Accuracy $\uparrow$ & $F_1$-Score $\uparrow$ & ROC-AUC $\uparrow$ & PR-AUC $\uparrow$ & Fit Time (s) \\
      \midrule
      \textbf{Ours: Decision Tree}  & \textbf{0.5036} & \textbf{0.2495} & \textbf{0.5992} & \textbf{0.9079} & \textbf{0.9131} & \textbf{0.9675} & \textbf{0.9671} & 4.53 s \\
      scikit-learn Decision Tree   & 0.5026 & 0.2493 & 0.6007 & 0.9078 & 0.9131 & 0.9674 & 0.9671 & 0.98 s \\
      \midrule
      \textbf{Ours: Pegasos SVM (RFF)} & --- & --- & --- & 0.8439 & 0.8512 & \textbf{0.9271} & \textbf{0.9375} & 10.90 s \\
      scikit-learn SGD Baseline    & --- & --- & --- & \textbf{0.8478} & \textbf{0.8578} & 0.9198 & 0.9255 & 0.63 s \\
      \midrule
      \textit{Bonus: Ours Ridge (Normal Eq.)} & 0.5172 & 0.3169 & 0.5773 & --- & --- & --- & --- & 0.09 s \\
      \textit{Bonus: Ours SVR (Sub-gradient)} & 0.5086 & 0.3132 & 0.5911 & --- & --- & --- & --- & 8.03 s \\
      \textit{Bonus: Ours Random Forest (46 trees)} & 0.4828 & --- & --- & --- & --- & --- & --- & --- \\
      \textit{Bonus: Ours Gradient Boosting (200 trees)} & \textbf{0.4552} & --- & --- & --- & --- & --- & --- & --- \\
      \bottomrule
    \end{tabular}%
    }
  \end{table*}

  \subsection{Do Our Models Match scikit-learn?}
  Table~\ref{tab:results} shows that our from-scratch decision tree gets very close to scikit-learn's optimized implementation. On Task A the RMSE differs by only $0.0010$ ($0.5036$ vs $0.5026$), and the $R^2$ values are almost the same ($0.5992$ vs $0.6007$). On Task B the results are practically the same: the same $F_1$ ($0.9131$), the same PR-AUC ($0.9671$), and almost the same ROC-AUC ($0.9675$ vs $0.9674$).

  On Task B our Pegasos SVM with RFF reaches an $F_1$ of $0.8512$. It ranks examples better than scikit-learn's linear SGD classifier (ROC-AUC $0.9271$ vs $0.9198$, PR-AUC $0.9375$ vs $0.9255$), which suggests that the non-linear Fourier features help. The SGD baseline is still slightly better on accuracy and $F_1$, though.

  \begin{figure}[ht]
    \centering
    \includegraphics[width=\linewidth]{figures/benchmark_tree.png}
    \caption{Comparison between our decision tree and scikit-learn on the Task A and Task B test splits.}
    \label{fig:benchmark_tree}
  \end{figure}

  \subsection{Cross-Validation}
  To check that our results do not depend on one lucky split, we ran 5-fold cross-validation on the training data ($62,922$ samples):
  \begin{itemize}
    \item \textbf{Regression RMSE}: Ours $0.2719 \pm 0.0083$, scikit-learn $0.2715 \pm 0.0086$.
    \item \textbf{Regression $R^2$}: Ours $0.8059 \pm 0.0085$, scikit-learn $0.8064 \pm 0.0088$.
    \item \textbf{Classification $F_1$}: Ours $0.9123 \pm 0.0016$, scikit-learn $0.9124 \pm 0.0013$.
    \item \textbf{Classification ROC-AUC}: Ours $0.9714 \pm 0.0012$, scikit-learn $0.9713 \pm 0.0010$.
  \end{itemize}
  The standard deviations are small (below 0.002 for classification), so performance does not change much between folds.

  \subsection{Hyperparameter Sensitivity and Bias--Variance Tradeoff}
  Figure~\ref{fig:tree_depth} shows how tree depth affects the model. For very deep trees ($d \ge 15$), training ROC-AUC goes up to $0.9996$ while validation ROC-AUC falls to $0.9252$, which is a clear sign of overfitting. Setting a minimum leaf size of $n_{\text{leaf}} = 10$ at $d = 12$ removes the overconfident leaves that contain a single sample, and raises validation ROC-AUC to $0.9660$.


  For the Pegasos SVM, the $(\lambda, \gamma)$ grid shows that a very large regularization ($\lambda \ge 0.1$) over-smooths the model: the weights shrink to almost zero ($\|\mathbf{w}\| \to 0$) and $F_1$ drops to about 0. The best balance was a small $\lambda$ with a moderate kernel scale ($\lambda = 10^{-4}$, $\gamma = 0.1$). It has a margin width of $2/\|\mathbf{w}\| = 0.1041$ and 23,840 margin support vectors, and gives the best validation $F_1$ of $0.8370$.

  \subsection{Error Analysis by Property Type}
  When we split the Task A test predictions by property type, the model works very differently across categories:
  \begin{itemize}
    \item \textbf{Apartments}: Predictions are accurate and consistent for \textit{Yeni tikili} ($n = 8,616$, RMSE $0.218$, $R^2 = 0.838$) and \textit{Köhnə tikili} ($n = 2,753$, RMSE $0.201$, $R^2 = 0.735$).
    \item \textbf{Commercial and land}: Errors are much larger for \textit{Obyekt} ($n = 623$, RMSE $1.256$, $R^2 = 0.062$) and \textit{Torpaq} ($n = 737$, RMSE $1.483$, $R^2 = 0.281$). Likely reasons are missing information such as street frontage, zoning permits, and access to utilities.
  \end{itemize}

  \begin{figure}[ht]
    \centering
    \includegraphics[width=\linewidth]{figures/tree_overfitting_depth.png}
    \caption{Bias-variance tradeoff for the depth sweep ($d \in [1, 18]$). Training ROC-AUC gets close to $0.9996$, while validation ROC-AUC drops for deep trees because the tree memorizes the training leaves.}
    \label{fig:tree_depth}
  \end{figure}

  \begin{figure}[ht]
    \centering
    \includegraphics[width=\linewidth]{figures/pr_curve.png}
    \caption{Precision-Recall curves on the test set for Task B. Both tree models keep high precision at high recall.}
    \label{fig:pr_curve}
  \end{figure}

  \begin{figure}[ht]
    \centering
    \includegraphics[width=\linewidth]{figures/tree_feature_importances.png}
    \caption{Top 15 features by Mean Decrease in Impurity (MDI). Living area, number of rooms, and the coordinates of premium Baku districts are the most important.}
    \label{fig:feature_importance}
  \end{figure}

  % =====================================================================
  \section{Discussion and Limitations}
  Our implementations perform close to the scikit-learn baselines, but they have some limits:
  \begin{enumerate}
    \item \textbf{Axis-aligned splits}: A tree can only split along one feature at a time, so it has trouble with diagonal patterns such as price gradients across Baku.
    \item \textbf{Limited random features}: Random Fourier Features allow non-linear boundaries, but $D = 256$ is limited by computing time, so very detailed local price patterns are only approximated.
    \item \textbf{Rare listings}: For luxury properties (above 500,000 AZN) and sparse rural regions (for example Astara and Gədəbəy), there are few examples, so prediction errors are larger.
  \end{enumerate}

  % =====================================================================
  \section{Conclusion}
  We built a recursive decision tree and a primal Pegasos SVM with Random Fourier Features from scratch, and tested them on real-estate data from \texttt{bina.az}. Our tree reaches results very close to scikit-learn ($F_1 = 0.9131$, and tree RMSE $0.5036$ vs $0.5026$). Cross-validation and the bias-variance experiments show that the results are stable. In the future we want to add spatial embeddings from geographic graph neural networks, and oblique (diagonal) splits in the tree.

  % =====================================================================
  \section*{Tools and Acknowledgements}
  We used AI coding assistants for template formatting, exploratory visualization scripts, and improving the documentation. The data comes from public listings on \texttt{bina.az}.

  % =====================================================================
  \begin{thebibliography}{9}
  \bibitem{pegasos}
  S.~Shalev-Shwartz, Y.~Singer, N.~Srebro, and A.~Cotter,
  ``Pegasos: primal estimated sub-gradient solver for SVM,''
  \emph{Mathematical Programming}, vol.~127, no.~1, pp.~3--30, 2011.

  \bibitem{breiman1984}
  L.~Breiman, J.~Friedman, C.~J.~Stone, and R.~A.~Olshen,
  \emph{Classification and Regression Trees}, CRC press, 1984.

  \bibitem{rahimi2007}
  A.~Rahimi and B.~Recht,
  ``Random features for large-scale kernel machines,''
  in \emph{Advances in Neural Information Processing Systems (NeurIPS)}, vol.~20, 2007.

  \bibitem{pedregosa2011}
  F.~Pedregosa et~al.,
  ``Scikit-learn: Machine learning in Python,''
  \emph{Journal of Machine Learning Research}, vol.~12, pp.~2825--2830, 2011.
  \end{thebibliography}

  \end{document}
